\clearpage
\InsertMark{ztex-4th}{}
\subsection{\titlepkg{thm} 库}\label{sec:theme-library-obsidian}
本 library 中定义了一系列的定理类主题以及环境图标(icon), 在加载 \pkg{theme} library 的同时，会自动导入 \pkg{tcolorbox}, \pkg{tikz} 和 \pkg{pifont} 
三个宏包. 同时也会加载 \pkg{tikz} 的 \pkg{fadings}, \pkg{calc} 两个库. 如此数量的宏包导入必然会拖慢整个文档的编译，请酌情考虑加载此 
library.

\vskip2em
\noindent\textcolor{red}{\sffamily NOTE: 
\begin{enumerate}
  \item 由于技术原因, 当用户需要加载 \pkg{thm} 库时, 必须将命令 \cs{zthmstyle}\texttt{\marg{style}} 置于 \cs{ztexloadlib}\zarg{thm} 之前;
  \item 若用户在自定义定理类环境样式时需要更改 \ztex{} 的默认配色, 请将 \cs{ztex_keys_set:nn} 或
    其它基于 \cs{keys_set:nn} 的命令放置于命令 \cs{zthmstylenew} 对应样式的 \meta{preamble} 中而非 \meta{option} 中, 
    否则 \ztex{} 中的一系列与 \cs{zcolorset} 相关的函数将失去对新定义数学类环境样式的色彩控制能力.
\end{enumerate}}


\clearpage
\begin{function}[updated=2025-04-25]{\zthmiconset}
  \begin{syntax}
    \cs{zthmiconset}\marg{key-value}
  \end{syntax}
  此命令用于设置定理类环境的图标, 仅能在导言区使用.
\end{function}

\begin{keyval}[parent=..]{axiom, definition, theorem, lemma, corollary, proposition, remark}
  \begin{syntax}
    axiom       = \meta{icon}>\dval{\ding{118}}
    definition  = \meta{icon}>\dval{\ding{168}} 
    theorem     = \meta{icon}>\dval{\(\heartsuit\)} 
    lemma       = \meta{icon}>\dval{\ding{68}}
    corollary   = \meta{icon}>\dval{\ding{168}} 
    proposition = \meta{icon}>\dval{\(\spadesuit\)} 
    remark      = \meta{icon}>\dval{\ding{102}}
    proof       = \meta{icon}>\dval{无}
    exercise    = \meta{icon}>\dval{无}
    example     = \meta{icon}>\dval{无}
    solution    = \meta{icon}>\dval{无}
    problem     = \meta{icon}>\dval{无}
  \end{syntax}
  上述键值配置为 \meta{style}\texttt{=paris} 时的样式, 其中 \meta{icon} 为一个合法的图标(文字).
\end{keyval}


一个基本的使用样例如下\zchcmd:
\begin{DocExample}*
\zthmiconset        
  {
    axiom       = \ding{118},
    definition  = \ding{168}, 
    theorem     = \(\heartsuit\), 
    lemma       = \ding{68},
    corollary   = \ding{168}, 
    proposition = \(\spadesuit\), 
    remark      = \ding{102},
  }
\end{DocExample}


\begin{function}[updated=2025-04-25]{\zthmiconuse}
  \begin{syntax}
    \cs{zthmiconuse}\marg{thm env name}
  \end{syntax}
  此命令用于使用定理类环境的图标, \meta{thm env name} 即为所有预定义的定理类环境名. 
  此命令在自定义定理环境样式时比较有用, 不推荐用户于正文中使用. \par
  一个基本的使用样例如下\zchcmd:
\end{function}
\begin{DocExample}*
\zthmiconuse{theorem}
\zthmiconuse{lemma}
\end{DocExample}


\begin{function}[updated=2025-04-25]{\zthmiconrm}
  \begin{syntax}
    \cs{zthmiconrm}
  \end{syntax}
  此命令会清除所有定理类环境的图标, 不推荐用户在正文中使用.
\end{function}



\begin{function}[updated=2024-12-05]{shadow}
  \begin{syntax}
    \cs{zthmstyle}\zarg{shadow}
  \end{syntax}
  加载此 library 后即可应用上述样式, 样式预览如下: 
\end{function}
\begingroup
\zthmstyle{shadow}
\begin{DocExample}*
% \ztexloadlib{alias}
\begin{remark}[thmstyle-shadow]
As any dedicated reader can clearly see, the Ideal of practical
reason is a representation of, as far as I know, the things in themselves; 
\begin{align}
\underset{}{\mathbf{v} \bigotimes \mathbf{w}} 
    & = \sum_{i=1}^3\left(a_{i1}u^iv^1+a_{i2}u^iv^2+a_{i3}u^iv^3\right) \\
    & = \int x \dd x = \frac12 x^2 + \R{C} 
  \end{align}  
As any dedicated reader can clearly see, the Ideal of practical
reason is a representation of, as far as I know, the things in themselves;%
\end{remark}
\end{DocExample}
\endgroup


\begin{function}[updated=2024-12-05]{paris}
  \begin{syntax}
    \cs{zthmstyle}\zarg{paris}
  \end{syntax}
  加载此 library 后即可应用上述样式, 样式预览如下: 
\end{function}
\begingroup
\zthmstyle{paris}
\zthmiconset
  {
    axiom       = \ding{118},
    definition  = \ding{168}, 
    theorem     = \(\heartsuit\), 
    lemma       = \ding{68},
    corollary   = \ding{168}, 
    proposition = \(\spadesuit\), 
    remark      = \ding{102},
  }
\begin{DocExample}*
% \ztexloadlib{alias}
\begin{axiom}[thmstyle-paris]
As any dedicated reader can clearly see, the Ideal of practical
reason is a representation of, as far as I know, the things in themselves; 
\begin{align}
\underset{}{\mathbf{v} \bigotimes \mathbf{w}} 
    & = \sum_{i=1}^3\left(a_{i1}u^iv^1+a_{i2}u^iv^2+a_{i3}u^iv^3\right) \\
    & = \int x \dd x = \frac12 x^2 + \R{C} 
  \end{align}  
As any dedicated reader can clearly see, the Ideal of practical
reason is a representation of, as far as I know, the things in themselves;%
\end{axiom}
\end{DocExample}
\endgroup


\begin{function}[updated=2024-12-05]{lapsis}
  \begin{syntax}
    \cs{zthmstyle}\zarg{lapsis}
  \end{syntax}
  加载此 library 后即可应用上述样式, 样式预览如下: 
\end{function}
\begingroup
\zthmtitleformat{\bfseries
\zthmname\ \zthmnumber
\zthmnotemptyTF{}{\\}
\zthmnote{}{}
}
\zthmiconset
  {
    axiom       = \ding{111},
    definition  = \ding{118}, 
    theorem     = \ding{169}, 
    lemma       = \ding{170},
    corollary   = \ding{168}, 
    proposition = \ding{125}, 
    remark      = \ding{46},
    proof       = , 
    exercise    = \ding{45},
    example     = ,
    solution    = \ding{45}, 
    problem     = ,
  }
\zthmstyle{lapsis}
\begin{DocExample}*
% \ztexloadlib{alias}
\begin{lemma}[thmstyle-lapsis]
As any dedicated reader can clearly see, the Ideal of practical
reason is a representation of, as far as I know, the things in themselves; 
\begin{align}
\underset{}{\mathbf{v} \bigotimes \mathbf{w}} 
    & = \sum_{i=1}^3\left(a_{i1}u^iv^1+a_{i2}u^iv^2+a_{i3}u^iv^3\right) \\
    & = \int x \dd x = \frac12 x^2 + \R{C} 
  \end{align}  
As any dedicated reader can clearly see, the Ideal of practical
\tcblower
\begin{align}
  \int x \dd x = \frac12 x^2 + \R{C}
\end{align}
reason is a representation of, as far as I know, the things in themselves;%
\end{lemma}
\end{DocExample}
\endgroup


\begin{function}[updated=2024-12-05]{elegant}
  \begin{syntax}
    \cs{zthmstyle}\zarg{elegant}
  \end{syntax}
  加载此 library 后即可应用上述样式, 样式预览如下: 
\end{function}
\begingroup
\zthmiconset{
  axiom       = \ding{118},
  definition  = \ding{168}, 
  theorem     = \(\heartsuit\), 
  lemma       = \ding{68},
  corollary   = \ding{168}, 
  proposition = \(\spadesuit\), 
  remark      = \ding{102},
}
\ExplSyntaxOn
\ztex_keys_set:nn {color}{
  axiom       = {HTML}{2c3e50},
  definition  = {RGB}{0, 166, 82},
  theorem     = {RGB}{255, 134, 23},
  lemma       = {RGB}{255, 134, 23},
  corollary   = {RGB}{255, 134, 23},
  proposition = {RGB}{0, 173, 247},
}
\ExplSyntaxOff
\zthmstyle{elegant}
\begin{DocExample}*
% \ztexloadlib{alias}
\begin{definition}[thmstyle-elegant]
As any dedicated reader can clearly see, the Ideal of practical
reason is a representation of, as far as I know, the things in themselves; 
\begin{align}
\underset{}{\mathbf{v} \bigotimes \mathbf{w}} 
    & = \sum_{i=1}^3\left(a_{i1}u^iv^1+a_{i2}u^iv^2+a_{i3}u^iv^3\right) \\
    & = \int x \dd x = \frac12 x^2 + \R{C} 
  \end{align}  
As any dedicated reader can clearly see, the Ideal of practical
reason is a representation of, as far as I know, the things in themselves;%
\end{definition}
\end{DocExample}
\endgroup


\begin{function}[added=2025-06-29]{tcb}
  \begin{syntax}
    \cs{zthmstyle}\zarg{tcb}
  \end{syntax}
  加载此 library 后即可应用上述样式, 样式预览如下: 
\end{function}
\begingroup
\ExplSyntaxOn\makeatletter
\ztex_keys_set:nn {color}
  {
    axiom       = {HTML}{2c3e50},
    remark      = purple!55!black,
    definition  = orange!55!black,
    theorem     = blue!55!black,
    lemma       = green!55!black,
    corollary   = green!55!black,
    proposition = {RGB}{0, 173, 247},
  }
\makeatother\ExplSyntaxOff
\zthmstyle{tcb}
\begin{DocExample}*
% \ztexloadlib{alias}
\begin{theorem}[thmstyle-tcb]
As any dedicated reader can clearly see, the Ideal of practical
reason is a representation of, as far as I know, the things in themselves; 
\begin{align}
\underset{}{\mathbf{v} \bigotimes \mathbf{w}} 
    & = \sum_{i=1}^3\left(a_{i1}u^iv^1+a_{i2}u^iv^2+a_{i3}u^iv^3\right) \\
    & = \int x \dd x = \frac12 x^2 + \R{C} 
  \end{align}  
As any dedicated reader can clearly see, the Ideal of practical
reason is a representation of, as far as I know, the things in themselves;%
\end{theorem}
\end{DocExample}
\endgroup



\begin{function}[updated=2024-12-05]{obsidian}
  \begin{syntax}
    \cs{zthmstyle}\zarg{obsidian}
  \end{syntax}
  加载此 library 后即可应用上述样式, 样式预览如下: 
\end{function}
\begingroup
\ExplSyntaxOn\makeatletter
\zthmtitleformat*{
  \noindent\sffamily\bfseries\textcolor{\thm@tmp@color}{
    \__ztex_thm_icon_use:o {\thm@tmp@name}
    \ \zthmname{\,:\,}\zthmnumber
  }
}
\__ztex_thm_icon_set:n 
  {
    axiom       = \ding{111},
    definition  = \ding{118}, 
    theorem     = \ding{169}, 
    lemma       = \ding{170},
    corollary   = \ding{168}, 
    proposition = \ding{125}, 
    remark      = \ding{46},
    proof       = , 
    exercise    = \ding{45},
    example     = ,
    solution    = \ding{45}, 
    problem     = ,
  }
\makeatother\ExplSyntaxOff
\zthmstyle{obsidian}
\begin{DocExample}*
% \ztexloadlib{alias}
\begin{proposition}[thmstyle-obsidian]
As any dedicated reader can clearly see, the Ideal of practical
reason is a representation of, as far as I know, the things in themselves; 
\begin{align}
\underset{}{\mathbf{v} \bigotimes \mathbf{w}} 
    & = \sum_{i=1}^3\left(a_{i1}u^iv^1+a_{i2}u^iv^2+a_{i3}u^iv^3\right) \\
    & = \int x \dd x = \frac12 x^2 + \R{C} 
  \end{align}  
As any dedicated reader can clearly see, the Ideal of practical
reason is a representation of, as far as I know, the things in themselves;%
\end{proposition}
\end{DocExample}
\endgroup